\documentclass[11pt,a4paper]{article}

\usepackage[utf8]{inputenc}
\usepackage[T1]{fontenc}
\usepackage{amsmath,amssymb,amsthm}
\usepackage{graphicx}
\usepackage{hyperref}
\usepackage{cleveref}
\usepackage{algorithm}
\usepackage{algorithmic}
\usepackage{booktabs}
\usepackage{listings}
\usepackage{geometry}
\geometry{margin=1in}

% Custom commands
\newcommand{\R}{\mathbb{R}}
\newcommand{\E}{\mathbb{E}}
\newcommand{\argmax}{\operatorname{argmax}}
\newcommand{\norm}[1]{\left\|#1\right\|}

% Theorem environments
\theoremstyle{definition}
\newtheorem{definition}{Definition}
\newtheorem{proposition}{Proposition}
\newtheorem{remark}{Remark}

\title{Action Memory in Multi-Agent Reinforcement Learning:\\Preventing Bang-Bang Control Through State Augmentation}

\author{Implementation Report\\CaNS Opposition Control with Deep Reinforcement Learning}

\date{\today}

\begin{document}

\maketitle

\begin{abstract}
This report presents the implementation of action memory in multi-agent deep reinforcement learning for turbulent flow control. We address the bang-bang control problem---characterized by high-frequency oscillations in control actions---through state augmentation with previous action information. We provide a rigorous mathematical formulation, explain the mechanism by which this prevents bang-bang behavior, and address the apparent redundancy in critic inputs. The implementation maintains full backwards compatibility while enabling smooth, adaptive temporal control.
\end{abstract}

\tableofcontents
\newpage

\section{Introduction}

\subsection{Problem Statement}

In multi-agent reinforcement learning for flow control, we observed undesirable bang-bang behavior: the policy exhibits high-frequency oscillations, rapidly switching between extreme action values (e.g., from $+1$ to $-1$). This behavior manifests as:

\begin{enumerate}
    \item High temporal gradient: $\norm{\nabla_t a}_{\text{RMS}} \approx 0.4$
    \item Excessive high-frequency content in FFT spectrum
    \item Suboptimal flow control due to actuator limitations and flow inertia
\end{enumerate}

\subsection{Root Cause Analysis}

\textbf{Paradox}: Bang-bang control achieves high drag reduction, yet is undesirable. Why?

The bang-bang behavior creates a \textbf{closed-loop feedback instability}:

\begin{enumerate}
    \item \textbf{Wave excitation}: Bang-bang actuation triggers high-frequency wave dynamics in the turbulent flow
    \item \textbf{Observation coupling}: These waves appear in velocity measurements $(\mathbf{u}, \mathbf{w})$
    \item \textbf{Reactive policy}: Policy $\pi(\mathbf{u}, \mathbf{w})$ responds to oscillating observations with oscillating actions
    \item \textbf{Feedback loop}: Actions reinforce wave dynamics, creating self-sustaining oscillations
    \item \textbf{Memory-less amplification}: Without action history, policy cannot detect or break this cycle
\end{enumerate}

The problem is NOT poor performance---drag reduction may be excellent. The issues are:
\begin{itemize}
    \item Actuators cannot physically achieve high-frequency switching
    \item May exploit numerical artifacts specific to the simulation
    \item Poor robustness to model mismatch or real-world conditions
    \item Unsustainable feedback dynamics
\end{itemize}

Mathematically, at time $t$ the policy receives state $\mathbf{s}_t = [\mathbf{u}_t, \mathbf{w}_t]$ and outputs action $\mathbf{a}_t = \pi(\mathbf{s}_t)$. At time $t+1$, it has no information about $\mathbf{a}_t$, so it cannot detect the oscillation pattern and learn to dampen it.

\subsection{Proposed Solution}

We augment the state representation to include the previous action:
\begin{equation}
    \mathbf{s}_t = [\mathbf{u}_t, \mathbf{w}_t, \mathbf{a}_{t-1}]
\end{equation}

This simple modification enables the policy to learn temporal smoothness adaptively while maintaining the Markov property.

\section{Mathematical Formulation}

\subsection{Standard MDP Without Action Memory}

The original formulation uses a Markov Decision Process (MDP) $\mathcal{M} = (\mathcal{S}, \mathcal{A}, P, R, \gamma)$ where:

\begin{align}
    \mathcal{S} &= \R^{2 \times 8 \times 8} \quad \text{(state space: velocity fields)} \\
    \mathcal{A} &= [-1, 1]^{8 \times 8} \quad \text{(action space: control field)} \\
    P(\mathbf{s}_{t+1} | \mathbf{s}_t, \mathbf{a}_t) &: \mathcal{S} \times \mathcal{A} \to \mathcal{S} \quad \text{(dynamics)} \\
    R(\mathbf{s}_t, \mathbf{a}_t) &: \mathcal{S} \times \mathcal{A} \to \R \quad \text{(reward)}
\end{align}

State representation for each agent $i$ at patch $(m,n)$:
\begin{equation}
    \mathbf{s}_t^{(i)} = \begin{bmatrix}
        \mathbf{u}_t^{(i)} \\
        \mathbf{w}_t^{(i)}
    \end{bmatrix} \in \R^{2 \times 8 \times 8}
\end{equation}

The policy network learns:
\begin{equation}
    \mathbf{a}_t^{(i)} = \pi_\theta(\mathbf{s}_t^{(i)}) = \pi_\theta(\mathbf{u}_t^{(i)}, \mathbf{w}_t^{(i)})
\end{equation}

\textbf{Issue}: This is memoryless. The policy at time $t+1$ has no knowledge of $\mathbf{a}_t^{(i)}$.

\subsection{MDP With Action Memory}

We augment the state space to include previous action:
\begin{equation}
    \mathcal{S}_{\text{aug}} = \R^{3 \times 8 \times 8} \quad \text{(augmented state space)}
\end{equation}

Augmented state representation:
\begin{equation}
    \mathbf{s}_t^{(i)} = \begin{bmatrix}
        \mathbf{u}_t^{(i)} \\
        \mathbf{w}_t^{(i)} \\
        \mathbf{a}_{t-1}^{(i)}
    \end{bmatrix} \in \R^{3 \times 8 \times 8}
\end{equation}

The transition dynamics become:
\begin{equation}
    P(\mathbf{s}_{t+1} | \mathbf{s}_t, \mathbf{a}_t) = P([\mathbf{u}_{t+1}, \mathbf{w}_{t+1}, \mathbf{a}_t] | [\mathbf{u}_t, \mathbf{w}_t, \mathbf{a}_{t-1}], \mathbf{a}_t)
\end{equation}

\begin{proposition}[Markov Property Preservation]
The augmented state representation $\mathbf{s}_t = [\mathbf{u}_t, \mathbf{w}_t, \mathbf{a}_{t-1}]$ preserves the Markov property:
\begin{equation}
    P(\mathbf{s}_{t+1} | \mathbf{s}_t, \mathbf{a}_t) = P(\mathbf{s}_{t+1} | \mathbf{s}_0, \mathbf{a}_0, \ldots, \mathbf{s}_t, \mathbf{a}_t)
\end{equation}
\end{proposition}

\begin{proof}
The next state depends only on current flow state $(\mathbf{u}_t, \mathbf{w}_t)$ and current action $\mathbf{a}_t$ through the Navier-Stokes equations. The augmented state includes $\mathbf{a}_{t-1}$, which becomes part of the next state representation. Since:
\begin{equation}
    \mathbf{s}_{t+1} = [\mathbf{u}_{t+1}, \mathbf{w}_{t+1}, \mathbf{a}_t]
\end{equation}
and $\mathbf{a}_t$ is determined by the current policy action, the Markov property holds.
\end{proof}

\subsection{Policy and Value Functions}

The policy network with action memory:
\begin{equation}
    \pi_\theta: \R^{3 \times 8 \times 8} \to [-1,1]^{8 \times 8}, \quad \mathbf{a}_t = \pi_\theta(\mathbf{u}_t, \mathbf{w}_t, \mathbf{a}_{t-1})
\end{equation}

The state-action value function (Q-function):
\begin{equation}
    Q^\pi(\mathbf{s}_t, \mathbf{a}_t) = \E_{\tau \sim \pi}\left[\sum_{k=0}^\infty \gamma^k r_{t+k} \Big| \mathbf{s}_t = [\mathbf{u}_t, \mathbf{w}_t, \mathbf{a}_{t-1}], \mathbf{a}_t \right]
\end{equation}

Bellman equation with action memory:
\begin{equation}
    Q(\mathbf{s}_t, \mathbf{a}_t) = r_t + \gamma \E_{\mathbf{s}_{t+1}, \mathbf{a}_{t+1}}\left[Q(\mathbf{s}_{t+1}, \mathbf{a}_{t+1})\right]
\end{equation}

Expanding:
\begin{equation}
    Q([\mathbf{u}_t, \mathbf{w}_t, \mathbf{a}_{t-1}], \mathbf{a}_t) = r_t + \gamma \E\left[Q([\mathbf{u}_{t+1}, \mathbf{w}_{t+1}, \mathbf{a}_t], \mathbf{a}_{t+1})\right]
\end{equation}

\textbf{Key observation}: The next state $\mathbf{s}_{t+1}$ includes the current action $\mathbf{a}_t$. This creates a temporal coupling.

\section{Mechanism: How Action Memory Prevents Bang-Bang}

\subsection{Temporal Coupling in Value Function}

\begin{proposition}[Temporal Coupling]
With action memory, the current action choice $\mathbf{a}_t$ directly affects the future state representation, creating an incentive for temporal consistency.
\end{proposition}

\begin{proof}
Consider two scenarios at time $t$ with identical flow states but different action histories:
\begin{align}
    \text{Scenario A:} \quad &\mathbf{s}_t^A = [\mathbf{u}_t, \mathbf{w}_t, +0.9] \\
    \text{Scenario B:} \quad &\mathbf{s}_t^B = [\mathbf{u}_t, \mathbf{w}_t, -0.7]
\end{align}

For a given current action $\mathbf{a}_t$, the next states differ:
\begin{align}
    \mathbf{s}_{t+1}^A &= [\mathbf{u}_{t+1}, \mathbf{w}_{t+1}, \mathbf{a}_t] \quad \text{if starting from A} \\
    \mathbf{s}_{t+1}^B &= [\mathbf{u}_{t+1}, \mathbf{w}_{t+1}, \mathbf{a}_t] \quad \text{if starting from B}
\end{align}

While the flow evolution $(\mathbf{u}_{t+1}, \mathbf{w}_{t+1})$ is identical, the Q-values differ:
\begin{equation}
    Q(\mathbf{s}_t^A, \mathbf{a}_t) \neq Q(\mathbf{s}_t^B, \mathbf{a}_t) \quad \text{in general}
\end{equation}

The Q-function learns to prefer actions that maintain consistency with $\mathbf{a}_{t-1}$ when beneficial.
\end{proof}

\subsection{Learning Smooth Control}

The policy learns smoothness through three mechanisms:

\subsubsection{Mechanism 1: Breaking Closed-Loop Feedback Instability}

\textbf{Critical observation}: Bang-bang control actually achieves \emph{high} drag reduction, not poor performance. The problem is more subtle:

\begin{enumerate}
    \item Bang-bang actuation triggers high-frequency wave dynamics in the flow
    \item These waves appear in velocity observations $(\mathbf{u}, \mathbf{w})$
    \item Policy sees oscillating observations $\to$ responds with bang-bang actions
    \item Actions reinforce wave dynamics $\to$ \textbf{closed-loop feedback instability}
\end{enumerate}

Mathematically, without action memory:
\begin{align}
    \text{Bang-bang } \mathbf{a}_t &\to \text{Wave dynamics in flow} \\
    \text{Waves in } (\mathbf{u}_{t+1}, \mathbf{w}_{t+1}) &\to \pi(\mathbf{u}_{t+1}, \mathbf{w}_{t+1}) = \text{Bang-bang } \mathbf{a}_{t+1} \\
    &\to \text{Self-reinforcing cycle}
\end{align}

The reward signal may actually favor bang-bang:
\begin{equation}
    \E[R_{\text{bang-bang}}] \geq \E[R_{\text{smooth}}] \quad \text{(bang-bang works well for drag!)}
\end{equation}

\textbf{However}, bang-bang is undesirable because:
\begin{enumerate}
    \item Not physically realizable (actuator bandwidth limitations)
    \item May exploit numerical artifacts or simulation resonances
    \item Poor robustness and generalization to real systems
    \item Creates unstable feedback loop with flow dynamics
\end{enumerate}

\textbf{With action memory}, the policy can detect and break this cycle:
\begin{equation}
    \pi(\mathbf{u}_t, \mathbf{w}_t, \mathbf{a}_{t-1}) \text{ learns: ``I'm oscillating } \to \text{ dampen the response''}
\end{equation}

The policy sees the full feedback loop and can learn to stabilize it, trading some drag reduction for smoother, more robust control.

\subsubsection{Mechanism 2: Smoothness Loss Regularization}

The actor loss includes spatial and consistency penalties:
\begin{equation}
    \mathcal{L}_{\text{actor}} = -\E_{\mathbf{s},\mathbf{a} \sim \pi}[Q(\mathbf{s}, \mathbf{a})] + \lambda_s \mathcal{L}_{\text{spatial}} + \lambda_c \mathcal{L}_{\text{consistency}}
\end{equation}

where:
\begin{align}
    \mathcal{L}_{\text{spatial}} &= \E\left[\norm{\nabla_x \mathbf{a}}^2 + \norm{\nabla_y \mathbf{a}}^2\right] \\
    \mathcal{L}_{\text{consistency}} &= \E\left[\sum_{\substack{i,j : \\\text{sim}(\phi_i, \phi_j) > \tau}} \max(0, \norm{\mathbf{a}_i - \mathbf{a}_j} - m)\right]
\end{align}

With action memory, these losses become more effective because the policy can plan ahead:
\begin{equation}
    \pi_\theta(\mathbf{u}_t, \mathbf{w}_t, \mathbf{a}_{t-1}) \text{ learns: ``Choose } \mathbf{a}_t \approx \mathbf{a}_{t-1} \text{ to minimize future smoothness loss''}
\end{equation}

\subsubsection{Mechanism 3: Value Function Smoothness}

The Q-function learns to be smooth in the action-history dimension:
\begin{equation}
    \frac{\partial Q([\mathbf{u}, \mathbf{w}, \mathbf{a}_{\text{prev}}], \mathbf{a}_{\text{curr}})}{\partial \mathbf{a}_{\text{curr}}} \text{ becomes coupled with } \mathbf{a}_{\text{prev}}
\end{equation}

This creates an implicit preference for:
\begin{equation}
    \mathbf{a}_{\text{curr}} \approx \mathbf{a}_{\text{prev}} \quad \text{when flow state is similar}
\end{equation}

\subsection{Closed-Loop Feedback Dynamics}

\begin{remark}[Bang-Bang as Closed-Loop Instability]
The bang-bang behavior can be modeled as a feedback system:
\begin{align}
    \text{Flow dynamics:} \quad &\mathbf{u}_{t+1} = f_{\text{NS}}(\mathbf{u}_t, \mathbf{a}_t) \\
    \text{Policy:} \quad &\mathbf{a}_t = \pi(\mathbf{u}_t) \\
    \text{Closed loop:} \quad &\mathbf{u}_{t+1} = f_{\text{NS}}(\mathbf{u}_t, \pi(\mathbf{u}_t))
\end{align}

Without action memory, if $\pi$ is highly reactive and flow dynamics $f_{\text{NS}}$ have resonances at certain frequencies, the closed loop can become unstable:
\begin{equation}
    \frac{\partial \pi}{\partial \mathbf{u}} \cdot \frac{\partial f_{\text{NS}}}{\partial \mathbf{a}} > 1 \quad \to \quad \text{Oscillatory instability}
\end{equation}

With action memory, the policy becomes:
\begin{equation}
    \mathbf{a}_t = \pi(\mathbf{u}_t, \mathbf{a}_{t-1})
\end{equation}

This adds \textbf{feedback damping}. The policy can learn:
\begin{equation}
    \frac{\partial \pi}{\partial \mathbf{a}_{t-1}} < 0 \quad \text{when } |\mathbf{a}_{t-1}| \text{ is large}
\end{equation}

effectively implementing a learned stabilizer that prevents runaway oscillations.
\end{remark}

\subsection{Adaptive Filtering vs Fixed Filtering}

\begin{remark}[Comparison with External Filtering]
External low-pass filtering applies a fixed rule:
\begin{equation}
    \mathbf{a}_t^{\text{filtered}} = \alpha \mathbf{a}_t^{\text{raw}} + (1-\alpha) \mathbf{a}_{t-1}^{\text{filtered}}
\end{equation}

This has two problems:
\begin{enumerate}
    \item Fixed $\alpha$ may be suboptimal (too smooth or too reactive)
    \item Credit assignment is broken: reward depends on $\mathbf{a}_t^{\text{filtered}}$ but policy outputs $\mathbf{a}_t^{\text{raw}}$
\end{enumerate}

With action memory, the policy learns adaptive filtering:
\begin{equation}
    \mathbf{a}_t = \pi_\theta(\mathbf{u}_t, \mathbf{w}_t, \mathbf{a}_{t-1}) = f_{\text{base}}(\mathbf{u}_t, \mathbf{w}_t) + \alpha_{\text{learned}}(\mathbf{u}_t, \mathbf{w}_t) \cdot \mathbf{a}_{t-1}
\end{equation}

where $\alpha_{\text{learned}}$ is adaptive and learned through experience.
\end{remark}

\section{Critic Architecture and Apparent Redundancy}

\subsection{Critic Input Composition}

The critic network receives:
\begin{equation}
    \text{Critic input: } [\mathbf{u}_t, \mathbf{w}_t, \mathbf{a}_{t-1}, \mathbf{a}_t] \in \R^{4 \times 64 \times 64}
\end{equation}

This appears redundant since:
\begin{equation}
    \mathbf{a}_{t-1}^{(t)} = \mathbf{a}_t^{(t-1)}
\end{equation}

\subsection{Why This Redundancy is Necessary}

\begin{proposition}[Critic Requires Full State and Action]
For proper Q-function approximation, the critic must receive both:
\begin{enumerate}
    \item The complete state: $\mathbf{s}_t = [\mathbf{u}_t, \mathbf{w}_t, \mathbf{a}_{t-1}]$
    \item The current action being evaluated: $\mathbf{a}_t$
\end{enumerate}
\end{proposition}

\begin{proof}
The Q-function is defined as:
\begin{equation}
    Q: \mathcal{S} \times \mathcal{A} \to \R
\end{equation}

With augmented state:
\begin{equation}
    Q(\mathbf{s}_t, \mathbf{a}_t) = Q([\mathbf{u}_t, \mathbf{w}_t, \mathbf{a}_{t-1}], \mathbf{a}_t)
\end{equation}

The critic must evaluate \emph{how good} is action $\mathbf{a}_t$ \emph{given} that the previous action was $\mathbf{a}_{t-1}$. Both pieces of information are essential:
\begin{itemize}
    \item $\mathbf{a}_{t-1}$ (from state): Historical context
    \item $\mathbf{a}_t$ (from action input): Current decision being evaluated
\end{itemize}
\end{proof}

\subsection{Information Encoded in Redundancy}

The apparent redundancy actually encodes valuable information:

\subsubsection{Action Change Magnitude}
\begin{equation}
    \Delta \mathbf{a}_t = \mathbf{a}_t - \mathbf{a}_{t-1}
\end{equation}

The critic can learn:
\begin{equation}
    Q([\mathbf{u}, \mathbf{w}, \mathbf{a}_{\text{prev}}], \mathbf{a}_{\text{curr}}) = Q_{\text{base}}(\mathbf{u}, \mathbf{w}, \mathbf{a}_{\text{curr}}) - \lambda_Q \norm{\mathbf{a}_{\text{curr}} - \mathbf{a}_{\text{prev}}}^2
\end{equation}

\subsubsection{Context-Dependent Action Evaluation}

The same action $\mathbf{a}_t$ has different value depending on action history:
\begin{align}
    Q([\mathbf{u}, \mathbf{w}, +0.9], +0.8) &> Q([\mathbf{u}, \mathbf{w}, -0.7], +0.8) \\
    \text{(Smooth continuation)} & \quad \text{(Large jump)}
\end{align}

\subsection{Critic Network Architecture}

The critic uses circular convolutions to process the full field:
\begin{equation}
    Q_\phi(\mathbf{s}_t, \mathbf{a}_t) = \text{MLP}\left(\text{Flatten}\left(\text{Conv}_{\text{circ}}\left(\begin{bmatrix}
        \mathbf{u}_t \\
        \mathbf{w}_t \\
        \mathbf{a}_{t-1} \\
        \mathbf{a}_t
    \end{bmatrix}\right)\right)\right)
\end{equation}

Input channels:
\begin{itemize}
    \item Without action memory: $[u, w, a_{\text{curr}}]$ — 3 channels
    \item With action memory: $[u, w, a_{\text{prev}}, a_{\text{curr}}]$ — 4 channels
\end{itemize}

\section{Implementation Details}

\subsection{Environment Modifications}

\subsubsection{State Representation}

Each agent observes an $8 \times 8$ patch:
\begin{equation}
    \text{obs}^{(i)} = \begin{cases}
        [\mathbf{u}^{(i)}, \mathbf{w}^{(i)}] & \text{without action memory (2 channels)} \\
        [\mathbf{u}^{(i)}, \mathbf{w}^{(i)}, \alpha \cdot \mathbf{a}_{t-1}^{(i)}] & \text{with action memory (3 channels)}
    \end{cases}
\end{equation}

where $\alpha$ is a scaling factor (default: 1.0).

\subsubsection{Previous Action Field Tracking}

The environment maintains:
\begin{equation}
    \mathbf{a}_{\text{field}}^{\text{prev}} \in \R^{64 \times 64}
\end{equation}

Updated after each step:
\begin{equation}
    \mathbf{a}_{\text{field}}^{\text{prev}} \leftarrow \mathbf{a}_{\text{field}}^{\text{curr}}
\end{equation}

The zero-mean constraint is applied in the training script \emph{before} actions reach the environment (see Section~7 for details).

\subsubsection{Episode Reset}
\begin{equation}
    \mathbf{a}_{\text{field}}^{\text{prev}} \leftarrow \mathbf{0}_{64 \times 64} \quad \text{at episode start}
\end{equation}

\subsection{Actor Network Architecture}

CNN encoder-decoder with variable input channels:
\begin{align}
    \text{Encoder:} \quad &\mathbf{h} = f_{\text{enc}}(\text{obs}) \in \R^{C \times 8 \times 8} \\
    \text{Similarity:} \quad &\boldsymbol{\phi} = \text{GAP}(g_{\text{sim}}(\mathbf{h})) \in \R^{d_{\text{sim}}} \\
    \text{Decoder:} \quad &\mathbf{a} = \tanh(f_{\text{dec}}(\mathbf{h})) \in [-1,1]^{8 \times 8}
\end{align}

Number of input channels:
\begin{equation}
    C_{\text{in}} = \begin{cases}
        2 & \text{without action memory} \\
        3 & \text{with action memory}
    \end{cases}
\end{equation}

\subsection{Critic Network Architecture}

Input preparation:
\begin{equation}
    \mathbf{X}_{\text{critic}} = \text{Reconstruct}(\{\text{obs}^{(i)}\}_{i=1}^{64}, \{\mathbf{a}^{(i)}\}_{i=1}^{64}) \in \R^{C_{\text{obs}} + 1 \times 64 \times 64}
\end{equation}

where $C_{\text{obs}} = 2$ or $3$ depending on whether action memory is enabled.

Forward pass:
\begin{align}
    \mathbf{h}_1 &= \text{CircularConv}(\mathbf{X}_{\text{critic}}) \\
    \mathbf{h}_2 &= \text{MaxPool}(\text{ReLU}(\text{GroupNorm}(\mathbf{h}_1))) \\
    &\vdots \\
    Q &= \text{Linear}(\text{Flatten}(\mathbf{h}_L)) \in \R
\end{align}

\subsection{Training Algorithm}

\begin{algorithm}
\caption{MADDPG with Action Memory}
\begin{algorithmic}[1]
\STATE Initialize actor $\pi_\theta$, critic $Q_\phi$, target networks $\pi_{\theta'}$, $Q_{\phi'}$
\STATE Initialize replay buffer $\mathcal{D}$
\FOR{episode $= 1, \ldots, N_{\text{episodes}}$}
    \STATE Reset environment, set $\mathbf{a}_{-1} = \mathbf{0}$
    \FOR{step $= 0, \ldots, T-1$}
        \STATE Observe $\mathbf{s}_t = [\mathbf{u}_t, \mathbf{w}_t, \mathbf{a}_{t-1}]$
        \STATE Select action $\mathbf{a}_t = \pi_\theta(\mathbf{s}_t) + \mathcal{N}(0, \sigma^2)$
        \STATE Execute $\mathbf{a}_t$, observe $r_t$, $\mathbf{u}_{t+1}$, $\mathbf{w}_{t+1}$
        \STATE Store transition $(\mathbf{s}_t, \mathbf{a}_t, r_t, \mathbf{s}_{t+1})$ in $\mathcal{D}$
        \IF{training\_step mod train\_freq $= 0$}
            \FOR{$g = 1, \ldots, G$}
                \STATE Sample minibatch $\{(\mathbf{s}_j, \mathbf{a}_j, r_j, \mathbf{s}_{j+1})\}$ from $\mathcal{D}$
                \STATE \textbf{Critic update:}
                \STATE \quad Compute target: $y_j = r_j + \gamma Q_{\phi'}(\mathbf{s}_{j+1}, \pi_{\theta'}(\mathbf{s}_{j+1}))$
                \STATE \quad Update: $\phi \leftarrow \phi - \eta_Q \nabla_\phi \frac{1}{B}\sum_j (Q_\phi(\mathbf{s}_j, \mathbf{a}_j) - y_j)^2$
                \STATE \textbf{Actor update:}
                \STATE \quad $\mathcal{L}_{\text{actor}} = -\frac{1}{B}\sum_j Q_\phi(\mathbf{s}_j, \pi_\theta(\mathbf{s}_j)) + \lambda_s \mathcal{L}_{\text{spatial}} + \lambda_c \mathcal{L}_{\text{consistency}}$
                \STATE \quad Update: $\theta \leftarrow \theta - \eta_\pi \nabla_\theta \mathcal{L}_{\text{actor}}$
                \STATE \textbf{Target update:}
                \STATE \quad $\theta' \leftarrow \tau \theta + (1-\tau)\theta'$
                \STATE \quad $\phi' \leftarrow \tau \phi + (1-\tau)\phi'$
            \ENDFOR
        \ENDIF
    \ENDFOR
\ENDFOR
\end{algorithmic}
\end{algorithm}

\section{Backwards Compatibility}

\subsection{Configuration-Based Selection}

Action memory is enabled via configuration:
\begin{verbatim}
observation:
  include_prev_action: true   # Enable (default: false)
  prev_action_scale: 1.0      # Scaling factor
\end{verbatim}

\subsection{Automatic Channel Detection}

The system automatically adapts:
\begin{align}
    C_{\text{in}}^{\text{actor}} &= \text{obs\_shape}[0] \in \{2, 3\} \\
    C_{\text{in}}^{\text{critic}} &= C_{\text{in}}^{\text{actor}} + 1 \in \{3, 4\}
\end{align}

\subsection{Network Initialization}

\begin{table}[h]
\centering
\begin{tabular}{lcc}
\toprule
Configuration & Actor Input & Critic Input \\
\midrule
Without action memory & $[u, w]$ (2 ch) & $[u, w, a_{\text{curr}}]$ (3 ch) \\
With action memory & $[u, w, a_{\text{prev}}]$ (3 ch) & $[u, w, a_{\text{prev}}, a_{\text{curr}}]$ (4 ch) \\
\bottomrule
\end{tabular}
\caption{Channel configuration for different modes}
\end{table}

\section{Expected Results}

\subsection{Smoothness Metrics}

\begin{table}[h]
\centering
\begin{tabular}{lcc}
\toprule
Metric & Without Memory & With Memory (Expected) \\
\midrule
Temporal gradient RMS & $\approx 0.4$ & $0.1$--$0.2$ \\
High-freq FFT content & High & Reduced \\
Action change rate & $> 50\%$ steps & $< 20\%$ steps \\
\bottomrule
\end{tabular}
\caption{Expected smoothness improvements}
\end{table}

\subsection{Performance Metrics}

\textbf{Important}: Drag reduction may slightly \emph{decrease} compared to bang-bang:
\begin{equation}
    \E[R_{\text{smooth}}] \approx \E[R_{\text{bang-bang}}] \quad \text{or slightly lower}
\end{equation}

This is an \textbf{acceptable trade-off} because:
\begin{enumerate}
    \item Bang-bang is not physically realizable (actuator limitations)
    \item Smooth control is more robust to model mismatch
    \item Avoids exploiting simulation-specific resonances
    \item Breaks unstable feedback loop with wave dynamics
\end{enumerate}

The goal is \textbf{not} to maximize drag reduction at all costs, but to find \emph{smooth, robust, physically realizable} control that still achieves good performance.

Total objective balances drag reduction with smoothness:
\begin{equation}
    J = \E[R_{\text{drag}}] - \lambda_s \E[\mathcal{L}_{\text{spatial}}] - \lambda_c \E[\mathcal{L}_{\text{consistency}}]
\end{equation}

Expected outcomes:
\begin{itemize}
    \item Drag reduction: $90\%$--$100\%$ of bang-bang performance (acceptable)
    \item Temporal smoothness: $4\times$ improvement (crucial)
    \item Physical realizability: Yes (vs no for bang-bang)
    \item Robustness: Improved (no feedback instability)
\end{itemize}

\subsection{Learning Dynamics}

Expected training progression:
\begin{enumerate}
    \item \textbf{Early training} (episodes 0--50): Policy ignores $\mathbf{a}_{t-1}$, still jumpy
    \item \textbf{Mid training} (episodes 50--150): Policy starts using $\mathbf{a}_{t-1}$ context
    \item \textbf{Late training} (episodes 150+): Smooth control emerges, temporal consistency learned
\end{enumerate}

\section{Zero-Mean Constraint: Critical Implementation Details}

\subsection{The Problem with Post-Processing Correction}

\textbf{Critical Discovery}: The standard approach of applying zero-mean correction as post-processing can \emph{itself cause} bang-bang behavior!

\subsubsection{How External Correction Creates Oscillations}

Standard implementation (PROBLEMATIC):
\begin{algorithm}
\caption{Post-Processing Zero-Mean (Problematic)}
\begin{algorithmic}[1]
\STATE Agents output raw actions: $\mathbf{a}_i^{\text{raw}}$ for $i = 1, \ldots, 64$
\STATE Collect into field: $\mathbf{A}^{\text{raw}} \in \R^{64 \times 64}$
\STATE Compute mean: $\mu = \frac{1}{N}\sum_{i,j} A^{\text{raw}}_{ij}$
\STATE Apply correction: $\mathbf{A}^{\text{exec}} = \mathbf{A}^{\text{raw}} - \mu$
\STATE Execute $\mathbf{A}^{\text{exec}}$, store in $\mathbf{a}_{\text{prev}}$ for next step
\end{algorithmic}
\end{algorithm}

\textbf{The oscillation mechanism}:
\begin{align}
t=0: &\quad \mu_0 = +0.3 \quad \to \quad \mathbf{A}^{\text{exec}}_0 = \mathbf{A}^{\text{raw}}_0 - 0.3 \quad \text{(shift DOWN)} \\
t=1: &\quad \mu_1 = -0.2 \quad \to \quad \mathbf{A}^{\text{exec}}_1 = \mathbf{A}^{\text{raw}}_1 + 0.2 \quad \text{(shift UP)} \\
t=2: &\quad \mu_2 = +0.4 \quad \to \quad \mathbf{A}^{\text{exec}}_2 = \mathbf{A}^{\text{raw}}_2 - 0.4 \quad \text{(shift DOWN)}
\end{align}

Even if $\mathbf{A}^{\text{raw}}_t$ is smooth, \textbf{varying $\mu_t$ creates oscillating $\mathbf{A}^{\text{exec}}_t$}!

\subsection{Gradient Mismatch Problem}

With post-processing correction:
\begin{equation}
\frac{\partial \mathcal{L}}{\partial \theta} = \frac{\partial \mathcal{L}}{\partial \mathbf{A}^{\text{exec}}} \cdot \frac{\partial \mathbf{A}^{\text{exec}}}{\partial \theta}
\end{equation}

But gradient computation assumes:
\begin{equation}
\mathbf{A}^{\text{exec}} = \pi_\theta(\mathbf{s}) \quad \text{(direct control)}
\end{equation}

Reality with correction:
\begin{equation}
\mathbf{A}^{\text{exec}} = \pi_\theta(\mathbf{s}) - \frac{1}{N}\sum_{i=1}^N \pi_\theta(\mathbf{s}_i) \quad \text{(indirect, coupled)}
\end{equation}

The Jacobian is:
\begin{equation}
\frac{\partial A_i^{\text{exec}}}{\partial A_j^{\text{raw}}} = \begin{cases}
1 - \frac{1}{N} & i = j \\
-\frac{1}{N} & i \neq j
\end{cases}
\end{equation}

\textbf{Gradients don't see this coupling!} Policy learns as if it directly controls executed actions.

\subsection{Solution: Differentiable Zero-Mean Correction}

Make the correction part of the computational graph:

\begin{algorithm}
\caption{Differentiable Zero-Mean (Correct)}
\begin{algorithmic}[1]
\STATE Forward pass all agents: $\mathbf{A}^{\text{raw}} = \pi_\theta(\mathbf{S})$ \quad (from $\tanh \to [-1,1]$)
\STATE Compute mean (differentiable): $\mu = \text{mean}(\mathbf{A}^{\text{raw}})$
\STATE Apply correction (differentiable): $\mathbf{A} = \mathbf{A}^{\text{raw}} - \mu$
\STATE Compute loss: $\mathcal{L} = -Q(\mathbf{S}, \mathbf{A}) + \lambda_s \mathcal{L}_{\text{spatial}} + \cdots$
\STATE Backpropagate: $\nabla_\theta \mathcal{L}$ includes effect of zero-mean constraint
\STATE Update: $\theta \leftarrow \theta - \eta \nabla_\theta \mathcal{L}$
\end{algorithmic}
\end{algorithm}

\textbf{Key benefits}:
\begin{enumerate}
    \item Gradients properly account for zero-mean constraint
    \item Policy learns to output actions with near-zero mean \emph{before} correction
    \item Correction magnitude $|\mu|$ becomes smaller during training
    \item Bang-bang oscillations from varying corrections are prevented
\end{enumerate}

\subsubsection{Interaction with Exploration Noise and Clipping}

During training, exploration noise is added:
\begin{equation}
\mathbf{A}^{\text{noisy}} = \mathbf{A}^{\text{raw}} + \mathcal{N}(0, \sigma^2 I)
\end{equation}

where $\mathbf{A}^{\text{raw}} \in [-1, 1]$ (from $\tanh$ activation). Adding noise can push actions outside $[-1, 1]$, requiring clipping:
\begin{equation}
\mathbf{A}^{\text{clipped}} = \text{clip}(\mathbf{A}^{\text{noisy}}, -1, 1)
\end{equation}

\textbf{Critical issue}: Clipping breaks the zero-mean property! If $\mathbf{A}^{\text{noisy}}$ is zero-mean, $\mathbf{A}^{\text{clipped}}$ may not be.

\textbf{Solution}: Apply zero-mean correction \emph{after} clipping:
\begin{equation}
\mathbf{A}^{\text{final}} = \mathbf{A}^{\text{clipped}} - \text{mean}(\mathbf{A}^{\text{clipped}})
\end{equation}

\textbf{Important}: When noise is annealed to zero ($\sigma = 0$ in late training), clipping is unnecessary:
\begin{itemize}
    \item Actor outputs via $\tanh$ are already in $[-1, 1]$
    \item No noise added $\to$ actions remain in $[-1, 1]$
    \item Clipping step can be skipped
    \item Zero-mean correction still applied as safety net (ensures exact zero-mean)
\end{itemize}

Full procedure:
\begin{algorithm}
\caption{Differentiable Zero-Mean with Noise Handling}
\begin{algorithmic}[1]
\STATE Forward pass: $\mathbf{A}^{\text{raw}} = \pi_\theta(\mathbf{S})$ \quad ($\in [-1,1]$ from $\tanh$)
\IF{training mode with noise}
    \STATE Generate zero-mean noise: $\boldsymbol{\epsilon} \sim \mathcal{N}(0, \sigma^2 I)$, $\boldsymbol{\epsilon} \leftarrow \boldsymbol{\epsilon} - \text{mean}(\boldsymbol{\epsilon})$
    \STATE Add noise: $\mathbf{A}^{\text{noisy}} = \mathbf{A}^{\text{raw}} + \boldsymbol{\epsilon}$ \quad (may exceed $[-1,1]$)
    \STATE Clip: $\mathbf{A}^{\text{clipped}} = \text{clip}(\mathbf{A}^{\text{noisy}}, -1, 1)$ \quad (breaks zero-mean!)
    \STATE Guarantee zero-mean: $\mathbf{A} = \mathbf{A}^{\text{clipped}} - \text{mean}(\mathbf{A}^{\text{clipped}})$
\ELSE
    \STATE No clipping needed (already in $[-1,1]$)
    \STATE Guarantee zero-mean: $\mathbf{A} = \mathbf{A}^{\text{raw}} - \text{mean}(\mathbf{A}^{\text{raw}})$
\ENDIF
\STATE Execute $\mathbf{A}$ (guaranteed zero-mean)
\end{algorithmic}
\end{algorithm}

\subsection{Per-Tile vs Global Zero-Mean Constraint}

\subsubsection{Per-Tile Constraint is Physically Invalid}

\textbf{Previously used (WRONG)}:
\begin{equation}
\mathcal{L}_{\text{zero}} = \frac{1}{N_{\text{agents}}}\sum_{i=1}^{N_{\text{agents}}} \left(\text{mean}(\mathbf{a}_i)\right)^2
\end{equation}

This penalizes non-zero mean per $8 \times 8$ tile.

\textbf{Why this is invalid}: Different flow regions require different average actuation:
\begin{itemize}
    \item Upstream regions may need net blowing: $\text{mean}(\mathbf{a}_{\text{upstream}}) > 0$
    \item Downstream regions may need net suction: $\text{mean}(\mathbf{a}_{\text{downstream}}) < 0$
    \item Separated regions require strong local forcing
\end{itemize}

Forcing each tile to zero-mean is an \textbf{unphysical artificial constraint}!

\subsubsection{Global Zero-Mean is Required}

\textbf{Correct approach (mass conservation)}:
\begin{equation}
\mathcal{L}_{\text{global\_mean}} = \left(\frac{1}{64 \times 64}\sum_{i,j} A_{ij}\right)^2
\end{equation}

This enforces:
\begin{equation}
\text{mean}(\mathbf{A}_{\text{field}}) = 0 \quad \text{(incompressibility)}
\end{equation}

while allowing:
\begin{equation}
\text{mean}(\mathbf{a}_i) \neq 0 \quad \text{for individual tiles}
\end{equation}

\subsection{Implementation}

\subsubsection{In Actor Network}

During training with exploration noise:
\begin{lstlisting}[language=Python]
# Forward pass for all agents
actions_raw = actor(obs_batch)  # [batch, 64, 8, 8], in [-1, 1] from tanh

# Add zero-mean exploration noise
if noise_scale > 0:
    noise = np.random.normal(0, noise_scale, size=actions_raw.shape)
    noise = noise - noise.mean()  # Ensure noise is zero-mean
    actions = actions_raw + noise  # May exceed [-1, 1]
    actions = np.clip(actions, -1, 1)  # Clipping needed
else:
    actions = actions_raw  # No clipping needed (already in [-1, 1])

# CRITICAL: Guarantee exact zero-mean after clipping
# (clipping can break zero-mean property)
actions = actions - actions.mean()

# Actions are now:
# 1. Exactly zero-mean (guaranteed)
# 2. In [-1, 1] range
# 3. Gradients flow through correction (differentiable)
\end{lstlisting}

During inference without noise (late training or evaluation):
\begin{lstlisting}[language=Python]
# Forward pass
actions_raw = actor(obs)  # In [-1, 1] from tanh

# No noise, no clipping needed
# Only zero-mean correction for exact guarantee
global_mean = actions_raw.mean()
actions = actions_raw - global_mean

return actions  # Guaranteed zero-mean, in [-1, 1]
\end{lstlisting}

\textbf{Key points}:
\begin{itemize}
    \item Clipping only required when adding noise ($\sigma > 0$)
    \item Without noise, $\tanh$ output is naturally in $[-1, 1]$
    \item Final zero-mean correction always applied (cheap safety net)
    \item Ensures exact $\text{mean}(\mathbf{A}) = 0$ for solver stability
\end{itemize}

\subsubsection{In Loss Function}

Old approach (WRONG):
\begin{equation}
\mathcal{L}_{\text{actor}} = -Q + \lambda_{\text{spatial}} \mathcal{L}_{\text{spatial}} + \lambda_{\text{zero}} \underbrace{\sum_i (\text{mean}(\mathbf{a}_i))^2}_{\text{per-tile, invalid}}
\end{equation}

New approach (CORRECT):
\begin{equation}
\mathcal{L}_{\text{actor}} = -Q + \lambda_{\text{spatial}} \mathcal{L}_{\text{spatial}} + \lambda_{\text{global}} \underbrace{(\text{mean}(\mathbf{A}_{\text{field}}))^2}_{\text{global only, valid}}
\end{equation}

\subsection{Expected Impact}

With these fixes:
\begin{enumerate}
    \item \textbf{Reduced bang-bang from corrections}: Policy learns to output near-zero mean $\to$ small $|\mu|$ $\to$ stable corrections
    \item \textbf{Proper credit assignment}: Gradients include zero-mean constraint effect
    \item \textbf{Physical validity}: Tiles can have non-zero mean as required by flow physics
    \item \textbf{Training stability}: Network aware of constraint during learning
    \item \textbf{Exact zero-mean guarantee}: Final correction after clipping ensures $\text{mean}(\mathbf{A}) = 0$ exactly for solver stability
    \item \textbf{Efficient inference}: When noise is annealed ($\sigma \to 0$), clipping is unnecessary since $\tanh$ output is naturally in $[-1, 1]$
\end{enumerate}

\begin{table}[h]
\centering
\begin{tabular}{lcc}
\toprule
Approach & Gradients See Constraint & Physically Valid \\
\midrule
Post-processing + per-tile loss & No & No \\
Differentiable + per-tile loss & Yes & No \\
Post-processing + global loss & No & Yes \\
\textbf{Differentiable + global loss} & \textbf{Yes} & \textbf{Yes} \\
\bottomrule
\end{tabular}
\caption{Comparison of zero-mean constraint approaches}
\end{table}

\section{Conclusion}

\subsection{Summary}

We have shown that augmenting the state with previous action:
\begin{equation}
    \mathbf{s}_t = [\mathbf{u}_t, \mathbf{w}_t] \quad \to \quad \mathbf{s}_t = [\mathbf{u}_t, \mathbf{w}_t, \mathbf{a}_{t-1}]
\end{equation}

provides the policy with the necessary information to learn smooth temporal control while:
\begin{enumerate}
    \item Preserving the Markov property
    \item Enabling adaptive (learned) temporal filtering
    \item Maintaining clean credit assignment
    \item Requiring minimal code changes
    \item Ensuring full backwards compatibility
\end{enumerate}

\subsection{Key Advantages Over Alternatives}

\begin{table}[h]
\centering
\small
\begin{tabular}{lccc}
\toprule
Approach & Markovian & Credit Assignment & Adaptive \\
\midrule
Temporal loss on replay & No & Confused & Partial \\
External filtering & Yes & Broken & No \\
\textbf{Prev action in obs} & \textbf{Yes} & \textbf{Clean} & \textbf{Yes} \\
\bottomrule
\end{tabular}
\caption{Comparison of approaches to temporal smoothness}
\end{table}

\subsection{Future Work}

Potential extensions:
\begin{enumerate}
    \item Multiple-step action history: $\mathbf{s}_t = [\mathbf{u}_t, \mathbf{w}_t, \mathbf{a}_{t-1}, \mathbf{a}_{t-2}]$
    \item Learned scaling: $\alpha(\mathbf{u}_t, \mathbf{w}_t)$ instead of fixed $\alpha$
    \item Action delta encoding: $[\mathbf{u}_t, \mathbf{w}_t, \Delta \mathbf{a}_{t-1}]$ instead of absolute $\mathbf{a}_{t-1}$
\end{enumerate}

\end{document}
